\documentclass[11pt]{article}
\usepackage[utf8]{inputenc}
\usepackage[T1]{fontenc}
\usepackage{amsmath,amssymb,bm}
\usepackage{graphicx}
\usepackage{booktabs}
\usepackage{float}
\usepackage{caption}
\usepackage[margin=1in]{geometry}
\usepackage[numbers,sort&compress]{natbib}
\usepackage{xcolor}
\usepackage{hyperref}
\hypersetup{colorlinks=true,linkcolor=blue,citecolor=blue,urlcolor=blue}
\usepackage{authblk}

\title{\textbf{Sparse Counterfactual Attribution of Crop-Yield Anomalies\\
to Extreme-Weather Drivers: A United States Study (1990--2025)}}
\author[1]{Tran Dai Phong Lam\thanks{Corresponding author: phonglam2599@gmail.com}}
\author[1]{Thu Le}
\author[1]{Nguyen Quoc Hung}
\author[1]{Nguyen Trung Trinh}
\affil[1]{FPT University, Ho Chi Minh City, Vietnam}
\date{}

\begin{document}
\maketitle

\begin{abstract}
Most crop-yield machine-learning studies optimize average predictive accuracy, yet
the years that matter most for food security and risk management are the abnormal
ones: regional crop failures. Average-accuracy models rarely say \emph{why} a
particular crop-region-year collapsed, or \emph{which} weather extreme was
responsible. We propose Sparse Counterfactual Anomaly Attribution (SCAA), an
event-level explanation method that detects low-yield anomalies and attributes each
one to a small, feasible set of extreme-weather drivers. Daily NASA POWER weather
(1989--2025), harmonized across twelve U.S.\ states and aggregated over
hemisphere-correct, crop-specific growing-season windows, is merged with USDA NASS
yields for Barley, Canola, Oats and Wheat. A forward-time gradient-tree model
predicts yield with a test $R^2$ of 0.809 (2016--2021) and 0.808 on an extended
2019--2025 horizon. Each crop-region yield series is detrended to separate
technological drift from weather-driven variation; years whose detrended residual
falls below $-1$ standard deviation are flagged as anomalies (215 of 1{,}227
observations, 17.5\%). For each anomaly, a sparse, feasibility-constrained
counterfactual computes the minimal weather change that would restore predicted
yield to its trend, yielding (i) a dominant driver and (ii) a weather-recoverable
fraction of the yield gap. Across anomalies the median recoverable fraction is 0.86,
and 89\% are classified as weather-driven; the remaining 11\% are flagged as
not weather-explainable, a useful signal for non-meteorological causes. Attribution
is validated against known events: the 2012 drought anomalies are attributed to heat
and growing-degree-day extremes, and the severe 2021--2022 anomalies to dry-spell and
moisture extremes; anomaly counts peak in the documented drought years 2002, 2011,
2021 and 2022. SCAA turns a black-box yield model into an interpretable,
event-level attribution tool. Results are predictive associations, not causal effects.

\smallskip
\noindent\textbf{Keywords:} crop yield; extreme weather; counterfactual explanation;
anomaly detection; interpretable machine learning; United States.
\end{abstract}

\section{Introduction}
\label{sec:intro}
Crop-yield prediction is dominated by models that minimize average error over all
crop-region-years \citep{vanKlompenburg2020CropYieldReview,Jabed2024Comprehensive}.
For planning and risk management, however, the decisive years are the abnormal ones:
the regional failures that strain supply and reveal climate vulnerability. An
average-accuracy model rarely answers the two questions a planner actually asks about
such a year: \emph{why} did this crop-region collapse, and \emph{which} weather
extreme was responsible? Global feature-importance scores describe the model in
aggregate, not the individual event.

This study reframes weather--yield interpretation as \emph{event-level attribution}.
We propose Sparse Counterfactual Anomaly Attribution (SCAA): detect low-yield
anomalies, then attribute each to a small, feasible set of extreme-weather drivers
using a sparse counterfactual. The method answers, for a specific bad year, ``the
smallest weather change that would have made this year normal,'' and identifies the
weather indicators that change. The approach is organized around three questions:
\begin{itemize}
\item \textbf{RQ1.} After separating long-run trend from year-to-year variation, which
crop-region-years are genuine low-yield anomalies?
\item \textbf{RQ2.} What fraction of each anomaly is recoverable by feasible weather
change, and which extreme-weather indicator is its dominant driver?
\item \textbf{RQ3.} Do the attributions agree with documented drought and heat events?
\end{itemize}
The contribution is a method that ports sparse counterfactual explanation
\citep{Mothilal2020DiCE} from instance-level classification to event-level yield-anomaly
attribution, with a continuous weather-recoverable fraction that also flags anomalies
that weather cannot explain. We study Barley, Canola, Oats and Wheat across twelve
U.S.\ states, 1990--2025, on a source-harmonized daily-weather record. All results are
predictive associations and planning signals, not causal effects.

\section{Related Work}
\label{sec:related}
Machine learning for yield prediction has been reviewed extensively, with weather,
soil and historical yield as common inputs across classical and deep models
\citep{vanKlompenburg2020CropYieldReview,Meghraoui2024DLReview,Jabed2024Comprehensive};
because no single model dominates across crops and regions, comparative studies are
standard \citep{Iniyan2023Techniques,Mahesh2024Gradient}. Interpretability work warns
that black-box yield models can mislead climate and planning conclusions without
careful, domain-aware explanation \citep{Hu2023ExplainableYield,Mohan2025AIXAI}.
Counterfactual explanation---minimal, feasible input changes that alter a
prediction---is a standard post-hoc tool \citep{Mothilal2020DiCE}; we adapt its
feasibility and sparsity principles from instance classification to yield-anomaly
attribution. Climate-attribution studies compare observed outcomes against a
counterfactual without climate trends \citep{Lobell2011ClimateTrends}; SCAA applies an
analogous logic at the level of individual crop-region-year anomalies, while
emphasizing that attributions are model-internal and not causal.

\section{Data}
\label{sec:data}
The unit of analysis is a crop-region-year; the target is yield in t/ha. We study the
overlap winter and spring crops Barley, Canola, Oats and Wheat across twelve U.S.\
states (Colorado, Illinois, Iowa, Kansas, Minnesota, Montana, Nebraska, North Dakota,
Oklahoma, South Dakota, Texas, Washington). Yields are from USDA NASS Quick Stats,
converted from bushels (or pounds, for canola) per acre to t/ha with crop-specific
test weights. Daily weather (maximum and minimum temperature, corrected precipitation,
all-sky surface radiation) is from NASA POWER, used uniformly for all states so that no
cross-source measurement difference confounds the analysis. The merged frame contains
1{,}257 crop-region-year observations spanning 1990--2025.

\section{Method}
\label{sec:method}

\subsection{Hemisphere-correct growing-season features}
\label{sec:fe}
Daily records are aggregated into interpretable growing-season indicators
(Figure~\ref{fig:design}): seasonal rainfall totals and means, dry-day counts and
maximum dry-spell length, long dry-spell event counts, heavy-rain day counts and
rolling 1/3/7-day maxima, heat-day counts and consecutive-heat runs, heatwave events,
heat- and growing-degree days, hot-dry days, frost and cold days, frost events,
radiation summaries, and early/mid/late thirds of the window. Crucially, the
aggregation window is crop- and latitude-specific rather than a single fixed window:
U.S.\ winter wheat (southern and Pacific-Northwest states) is aggregated over
September--June, while spring crops (barley, oats, canola, and northern-plains spring
wheat) are aggregated over April--September. This phenology-correct windowing captures
the over-wintering period of winter wheat (e.g.\ the establishment-stage frost
exposure) that a fixed warm-season window omits.

\subsection{Detrending and anomaly definition}
\label{sec:detrend}
Long-run yield change is dominated by technological and managerial trend rather than by
weather, so we separate the two before attribution. For each crop-region series we fit
an ordinary-least-squares trend $\hat y_{c,r}(t)=\beta_0+\beta_1 t$ and form the
residual $e_{c,r,t}=y_{c,r,t}-\hat y_{c,r}(t)$, standardized within series to
$z_{c,r,t}=e_{c,r,t}/\sigma_{e}$. A crop-region-year is a \emph{low-yield anomaly} if
$z_{c,r,t}<-1$. This removes the trend that would otherwise dominate any attribution and
isolates the weather-driven shortfall.

\subsection{Yield model}
\label{sec:model}
An Extremely Randomized Trees regressor (400 trees, minimum leaf size 2) maps the
growing-season features, the year index, location and crop/region identifiers to yield.
Numeric features are median-imputed and standardized; categorical features are one-hot
encoded. The model is fit under a forward-time protocol (train $\le2015$, test
$\ge2016$).

\subsection{Sparse counterfactual anomaly attribution (SCAA)}
\label{sec:scaa}
Let $\hat f$ be the yield model and, for an anomaly $(c,r,t)$, let
$\tau=\hat y_{c,r}(t)$ be the trend-expected ``normal'' yield. With $\mathbf w$ the
observed weather vector, SCAA seeks the minimal feasible weather change that raises the
predicted yield toward $\tau$:
\begin{equation}
\bm\delta^\star=\arg\min_{\bm\delta}\;\|\bm\delta\oslash\bm\sigma_r\|_1
\quad\text{s.t.}\quad \hat f(\mathbf w+\bm\delta)\ \uparrow\ \tau,\;\;
\mathbf w+\bm\delta\in\mathcal W_r,
\end{equation}
where $\bm\sigma_r$ is the per-feature regional standard deviation and $\mathcal W_r$
the region's observed feasible box (preventing out-of-distribution counterfactuals).
Because $\hat f$ is a tree ensemble, we solve by greedy coordinate search over
interpretable indicators with a sparsity budget. Rather than a binary repaired/not, we
report the continuous \emph{weather-recoverable fraction}
\begin{equation}
\phi_{c,r,t}=\frac{\hat f(\mathbf w+\bm\delta^\star)-\hat f(\mathbf w)}{\tau-\hat f(\mathbf w)},
\end{equation}
the share of the modelled yield gap that feasible weather change can recover, and the
\emph{dominant driver}, the indicator with the largest standardized change in
$\bm\delta^\star$. An anomaly is \emph{weather-driven} if $\phi\ge0.5$; a low $\phi$
flags an anomaly that weather alone cannot explain. All attributions are post-hoc and do
not influence model fitting.

\section{Results}
\label{sec:results}

\subsection{Predictive performance}
\label{sec:perf}
On harmonized data with phenology-correct windows the yield model attains a test $R^2$
of 0.809 (RMSE 0.525 t/ha) over 2016--2021, and 0.808 (RMSE 0.536) over an extended
2019--2025 horizon, confirming stable forward-time skill into recent years
(Figure~\ref{fig:scatter}). Predictability is markedly regime-dependent: spring crops
reach $R^2=0.842$ while winter wheat, despite its correct September--June window,
remains harder at $R^2=0.543$, reflecting the larger weather-driven variance of
semi-arid winter-wheat regions.

\begin{figure}[h]
\centering
\includegraphics[width=0.92\linewidth]{figures/fig02_daily_extreme_feature_design.png}
\caption{From daily weather to interpretable growing-season indicators. Daily records are
aggregated over crop- and latitude-specific windows into extreme-weather features used
for prediction and anomaly attribution.}
\label{fig:design}
\end{figure}

\begin{figure}[h]
\centering
\includegraphics[width=0.5\linewidth]{figures/figUS_scatter.png}
\caption{Observed versus predicted yield, forward-time test 2016--2025. Calibration is
close to the identity line across crops.}
\label{fig:scatter}
\end{figure}

\subsection{Anomaly detection}
\label{sec:anom}
After detrending, 215 of 1{,}227 crop-region-years (17.5\%) are flagged as low-yield
anomalies ($z<-1$). Their temporal distribution (Figure~\ref{fig:timeline}) concentrates
in documented U.S.\ drought years: anomaly counts peak in 2002, 2011, 2021 and 2022, with
2012 also prominent---each a recognized drought or heat year on the Plains. This temporal
alignment is independent corroboration that the detrended anomalies capture genuine
weather-driven failures rather than statistical noise.

\begin{figure}[h]
\centering
\includegraphics[width=0.85\linewidth]{figures/figH_anomaly_timeline.png}
\caption{Low-yield anomalies per year; major drought years (2012, 2021, 2022) highlighted.
Anomaly frequency tracks documented drought episodes.}
\label{fig:timeline}
\end{figure}

\subsection{Counterfactual attribution}
\label{sec:attr}
SCAA assigns each anomaly a weather-recoverable fraction and a dominant driver
(Figure~\ref{fig:scaa}). The median recoverable fraction is 0.86, and 89\% of anomalies
are weather-driven ($\phi\ge0.5$); the remaining 11\% have low $\phi$ and are flagged as
not weather-explainable, a useful pointer to non-meteorological causes (pests,
management, market-driven abandonment, or moisture excess beyond the model's sensitivity).
The dominant drivers form an agronomically coherent set (Table~\ref{tab:drivers}): long
dry-spell events and frost events lead, followed by seasonal rainfall deficits, heat-day
exposure, excess-wet days and radiation, reflecting that U.S.\ crop failures arise from
both drought/heat and, less often, cold and excess-moisture extremes.

\begin{table}[h]
\centering
\caption{Most frequent dominant weather driver across weather-driven anomalies.}
\label{tab:drivers}
\small
\begin{tabular}{lr}
\toprule
Dominant driver & \# anomalies \\
\midrule
Long dry-spell events ($\ge$14 days) & 23 \\
Frost events ($\ge$2 consecutive days) & 18 \\
Seasonal rainfall (deficit) & 16 \\
Heat days ($T_{\max}\ge30^\circ$C) & 15 \\
Wet days (excess) & 12 \\
Seasonal radiation & 11 \\
Mean maximum temperature & 9 \\
Heatwave events (3-day, $30^\circ$C) & 9 \\
\bottomrule
\end{tabular}
\end{table}

\begin{figure}[h]
\centering
\includegraphics[width=0.95\linewidth]{figures/figG_us_methodA.png}
\caption{SCAA outputs. Left: distribution of the weather-recoverable fraction; 89\% of
anomalies exceed the 0.5 weather-driven threshold (median 0.86). Right: frequency of the
dominant weather driver across attributed anomalies.}
\label{fig:scaa}
\end{figure}

\subsection{Event validation}
\label{sec:valid}
The attributions agree with known events. In the 2012 Plains drought, the Kansas Oats
anomaly is attributed to growing-degree-day and heat extremes (recoverable fraction
0.93), North Dakota Canola to consecutive heat days (0.86), and South Dakota Barley to
heatwave events (0.60)---all heat-driven, consistent with the 2012 heat-drought. The most
severe recent anomalies, in 2021--2022, are attributed predominantly to dry-spell and
dry-day extremes (e.g.\ Oklahoma Oats 2022 and South Dakota Barley 2021), matching the
2021 Northern-Plains/Pacific-Northwest drought and the 2022 droughts. A minority of 2021
North Dakota anomalies have low recoverable fraction with excess-moisture drivers,
correctly flagged by SCAA as not cleanly weather-explainable.

\section{Discussion}
\label{sec:discussion}
SCAA converts a competitive but opaque yield model into an interpretable, event-level
tool. The two-stage design---detrend, then attribute the residual anomaly---matters: long-run
yield change is trend-dominated, so attributing raw change to weather would be
misleading, echoing cautions in explainable yield modelling \citep{Hu2023ExplainableYield}
and climate-trend attribution \citep{Lobell2011ClimateTrends}. The continuous
weather-recoverable fraction is more honest than a binary explanation: it quantifies how
much of each failure weather can account for and, by flagging low-$\phi$ anomalies,
distinguishes weather failures from those requiring non-weather explanations. Validation
against 2012 and 2021--2022 shows the attributions recover the correct physical driver for
documented events.

\section{Limitations and Future Work}
\label{sec:limits}
Attributions are model-internal predictive associations, not causal effects, and depend
on the model's weather sensitivity; the greedy counterfactual margin is a conservative
lower bound. State-level aggregation cannot resolve within-state management, irrigation,
cultivar or soil variation. Linear detrending may under- or over-fit non-linear technology
curves for some series. Future work should add irrigation and management covariates,
finer spatial resolution, bootstrap intervals on the recoverable fraction, and a synthetic
benchmark with known weather/trend contributions to validate the attribution itself.

\section{Conclusion}
\label{sec:conclusion}
We presented SCAA, a sparse-counterfactual method that detects low-yield anomalies and
attributes each to its dominant extreme-weather driver with a weather-recoverable
fraction. On twelve U.S.\ states and four crops (1990--2025), a forward-time model reaches
$R^2=0.81$, anomalies align with documented droughts, 89\% of anomalies are weather-driven,
and event-level attributions recover the correct drivers for the 2012 and 2021--2022
events while honestly flagging the cases weather cannot explain. SCAA makes
extreme-weather yield interpretation operational at the level of the individual events
that matter most for risk.

\small
\bibliographystyle{plainnat}
\bibliography{references}
\end{document}
